\begin{papercard}{hudig2025}
\cf{Stages} \stagetag{D} \stagetag{W} \stagetag{H} (also \stagetag{S}). A three-method study of fertility trackers: network data-flow monitoring of eight products, a contextual questionnaire ($n=162$), and four co-design focus groups ($n=15$). It yields ten transparency and control mechanisms organised by the \emph{product usage lifecycle}.

\cf{Aims} (1) Insight into fertility-tracker data sharing via transmissions; (2) nuanced understanding of user preferences on intimate data sharing; (3) actionable transparency and control mechanisms \pc{hudig2025}{2}{This paper has three key aims}. Definitions: \emph{transparency} enables understanding and informed action \pc{hudig2025}{4}{when realised effectively, en}; \emph{control} is the ability, not obligation, to manage, correct, delete or decide \pc{hudig2025}{4}{is the ability (yet not the obligation) to manage}. Fertility trackers are defined functionally as stand-alone apps or devices paired with apps for self-monitoring cycles, pregnancy and menopause \pc{hudig2025}{3}{are technical tools designed to help women self-monitor their reproductive cycles}.

\cf{Phase 1: data flows} Eight products pseudonymised: four apps, plus devices (wearables, thermometer, hormone monitor). Selection used popularity (Play downloads), diversity, an \emph{on-device benchmark app}, UK availability and Bluetooth \pc{hudig2025}{5}{As our goal was to gain a general}. Monitoring access point: a Raspberry Pi with tcpdump and a tablet on Android 10. Content was mostly encrypted, so only metadata could be analysed \pc{hudig2025}{5}{The data flow monitoring methodology and}. Four experiment phases with time-logged actions: \emph{setup}, \emph{interactions}, \emph{idle} (70 minutes), \emph{custom privacy settings} \pc{hudig2025}{5}{For each app, we conducted four phases of experiments}. A control baseline of 44 hours of tablet traffic without apps removed background IPs \pc{hudig2025}{5}{To isolate app data from other system}. IPs were mapped to organisations and countries via IPInfo plus manual research, with policies read where relevant.

\cf{Phase 2: questionnaire} Qualtrics. Scenario with a fictional company asking about willingness to share five data types with randomised stakeholders, then the same with \emph{de-identified} data. A 100-point allocation over ten GDPR-derived factors (transparency, purpose limitation, rights, security, de-identification) \pc{hudig2025}{6}{In (ii), respondents identified factors influencing their comfort}. Open to all adults including non-users and all genders, after 12 pilot responses. Analysis: descriptives and $\chi^2$, plus single-researcher content analysis of open answers \pc{hudig2025}{7}{A single researcher employed content analysis}.

\cf{Phase 3: co-design focus groups} Four online groups of 3--5 people (up to 90 minutes), led by a woman, each including one non-user to counter design fixation \pc{hudig2025}{7}{Given the sensitive nature of the topic, we curated small}. Structure: discussion, co-design on a lifecycle frame (\emph{selection, setup, use, post-use}) on Miro, and reflection on responsibilities \pc{hudig2025}{7}{product usage lifecycle stages: Selection of a fertility tracker}. Framework analysis by one researcher with team input, predefined themes plus emergent sub-codes, giving 40 concepts in ten mechanisms \pc{hudig2025}{7}{We analysed the focus group data using}. Ethics: optional sensitive items, consent, a safe-space protocol, a \pounds 20 voucher \pc{hudig2025}{8}{Given the nature of the topic, particular attention was paid to sen}.

\cf{Key findings} On average 13.5 organisations in 7.5 countries per tracker, and transmissions coincided with intimate inputs (cycle dates, intercourse, temperature, hormone tests) \pc{hudig2025}{9}{Our analysis revealed clear}. The on-device app (Cherry) showed no use-time transmissions, consistent with its claim. Privacy settings reduced flows in some apps but not visibly in others, e.g.\ an ``anonymous mode'' \pc{hudig2025}{10}{The other apps showed less obvious differences when their pri}. Willingness to share depends mainly on the \emph{stakeholder}: 88\,\% would share nothing with ad platforms and 65\,\% nothing with employers \pc{hudig2025}{11}{They were particularly unwilling to share their data with an Online}. Ten mechanisms (Table 6), e.g.\ store summaries, marked mandatory fields, toggles, selective agreement to terms, real-time sharing notices, a ``Your Data'' tab, export tools, post-use deletion information \pc{hudig2025}{17}{Simple Summaries on Download Platforms}. Participants feared \emph{user burden} and wanted regulators and providers to share responsibility.

\cf{Limitations stated} Eight products, metadata only, inferred organisations, imperfect isolation, a young educated mostly European sample, and lifecycle framing that may have shaped where ideas surfaced \pc{hudig2025}{18}{We highlight several limitations of our study}.

\cf{Use for this thesis} The \emph{lifecycle frame} (selection, setup, use, post-use) is directly usable to structure walkthrough tasks: what a user can learn before install, at onboarding, during tracking and after deletion. Mechanisms 4, 7, 9 and 10 give observable items such as marked mandatory fields, purpose icons, ``Your Data'' tab and post-deletion information. The stakeholder list (provider, research, health services, advertisers, insurers, partners, employers) is a literature-backed \emph{actor set} for harm-tree recipients. Traffic monitoring and co-design are out of scope. Use findings as context, not as observations about the German sample.
\end{papercard}
